% ════════════════════════════════════════
% 摘要（00.摘要.tex）
% ════════════════════════════════════════
\begin{abstract}

本文针对微网与外部电网的电力调控问题，结合\textbf{线性规划与滚动时域优化方法}，提出了系统的购电与储能联合调控方案。

\textbf{针对问题一:} 本问是\textbf{单日确定性优化}问题，需在电价、负载相同且光伏预测已知时制定储能 0:00 与 24:00 储电量相等的日计划购电策略。对附件一进行时间对齐，重构为 144 个 10 分钟间隔；以购电功率、充放电功率与储电量为决策变量，建立功率平衡、储能状态递推、容量与功率约束；以全天购电费用最小为目标，加入日循环约束并以附录1给定的 6000 kWh 固定 0:00 储电量，构建\textbf{线性规划模型}求解；得到全天购电量 \textbf{59482.70 kWh}、购电费 \textbf{35126.95 元}，储能 0:00 与 24:00 储电量均为 \textbf{6000 kWh}。结果表明，储能低价充电、高价放电可显著降低购电费用。

\textbf{针对问题二:} 本问是\textbf{全年滚动优化}问题，需在负载、光伏随时间变化且供电不足按 5 倍电价紧急购电时制定每天 0:00 的计划购电策略。读取全年 365 天的负载与光伏并按天切分；以 1 月 1 日 0:00 储电量 6000 kWh 为初值，逐日求解线性规划并传递日末储电量；建立紧急购电的 5 倍电价\textbf{惩罚机制}。其中 2025 年 1 月作为\textbf{预热期}，用于使储电量摆脱初始条件影响、进入稳定日循环；按附件 5 要求，结果文件与费用汇总均只统计 2 月 1 日—12 月 31 日（334 天），该口径下购电费用合计 \textbf{12254765.72 元}。结果表明，规划时已知真实负载与光伏，紧急购电量为零，储能稳定完成低充高放套利。

\textbf{针对问题三:} 本问是\textbf{滚动预报下的调整优化}问题，需在 0:00、6:00、12:00、18:00 利用更新预报调整购电并计及违约与超额费用。将小时光伏预报插值到 10 分钟；在 0:00 制定计划购电、在其余时刻滚动重优化；用上境变量将计划偏离的分段费用线性化，并以实际供电缺口核算紧急购电；得到（同样仅统计 2 月 1 日—12 月 31 日）含调整总费用 \textbf{14337999.15 元}，较仅用 0:00 计划降低 \textbf{444073.36 元}。结果表明，引入其他预报时刻的调整可使费用降低约 \textbf{3.00\%}。

\textbf{针对问题四:} 本问是\textbf{波动电价下的重算}问题，需在实时波动电价下重算问题二与问题三。读取全年波动电价数据；将模型电价参数逐\textbf{间隔替换}；重新求解逐日计划与滚动调整两类线性规划；得到波动电价下问题二购电费合计 \textbf{12841054.36 元}、问题三总费用 \textbf{14956526.18 元}。结果表明，波动电价增大了微网购电成本，储能套利仅能部分对冲电价波动的影响。

综上，本文完成了四个问题的建模与求解，模型全局最优、结论可靠，可为微网经济调度提供理论依据和决策参考。

\keywords{\textbf{微网}；\textbf{储能调度}；\textbf{线性规划}；\textbf{滚动时域优化}；\textbf{光伏预报}}

\end{abstract}
